\documentclass[11pt,a4paper]{article}
\usepackage{templates/modeling-notes}
\usepackage{needspace}
\hypersetup{pdftitle={Chapter 5 -- Sensitivity analysis and Lagrange multipliers}}
\title{\textbf{Mathematical Modeling}\\[2mm]
\large Chapter 5: Sensitivity Analysis and Lagrange Multipliers}
\author{Personal study notes}
\date{}
\begin{document}
\maketitle
{\small\tableofcontents}
\clearpage

\section{What an optimal basis tells us}\label{sec:basis}
Suppose a linear program has already been solved. Sensitivity analysis asks what
happens when its data change, and whether the final basis can be reused.
We work with
\[
  \max c^Tx,\qquad Ax=b,\qquad x\geq0,\qquad
  A\in\R^{m\times n},\quad\operatorname{rank}A=m.
\]
An ordered basis $B=(j_1,\ldots,j_m)$ selects a nonsingular matrix $A_B$.
Let $N$ be the remaining indices, and define
\begin{equation}\label{eq:basis-data}
  \bar A=A_B^{-1}A,\quad\bar b=A_B^{-1}b,\quad
  y=A_B^{-T}c_B,\quad q=A^Ty-c=\bar A^Tc_B-c.
\end{equation}
Our sign convention is $q_j=z_j-c_j$, the negative of the usual reduced cost.
The basic solution and objective dictionary are
\begin{equation}\label{eq:dictionary}
  x_N=0,\quad x_B=\bar b,\qquad
  c^Tx=z_0-q_N^Tx_N,\quad z_0=c_B^T\bar b=b^Ty.
\end{equation}
Thus $\bar b\geq0$ establishes primal feasibility, and $q\geq0$ establishes
dual feasibility for $\min b^Ty$ subject to $A^Ty\geq c$, $y$ free.
Together they certify optimality.

Here an \emph{optimal basis} means a basis passing both tests. An optimal point
can have several bases, and at a degenerate optimal point a particular basis may
fail $q\geq0$. The calculations below follow the chosen certifying basis.
Losing its certificate does not, by itself, show that an old degenerate point
has lost optimality.

\begin{center}
  \begin{tabular}{>{\raggedright\arraybackslash}p{2.5cm}>{\raggedright\arraybackslash}p{5.1cm}>{\raggedright\arraybackslash}p{5.5cm}}
    \toprule
    What is retained? & Meaning & What can still change?\\\midrule
    Basis & Same column indices; both tests pass & Basic values, objective, dual prices\\
    Point & Same decision vector remains optimal & Basis, objective value, multipliers\\
    Optimal value & Same best objective value & Optimizers and bases\\\bottomrule
  \end{tabular}
\end{center}
Keeping a basis is often the useful goal: it gives explicit formulas rather
than requiring a new solve. Simultaneous changes are analyzable;
they must be tested together, not against separate one-parameter intervals.

\begin{figure}[H]\centering
  \begin{tikzpicture}[font=\small]
    \node[draw=modelblue,rounded corners,align=center,text width=3cm] (data) at (0,0)
    {Changed model\\$A',b',c'$};
    \node[draw=modelblue,rounded corners,align=center,text width=4.4cm] (tests) at (4.8,0)
    {Recompute with the old basis\\$\bar b'=A_B'^{-1}b'$\\$q'=A'^TA_B'^{-T}c_B'-c'$};
    \node[draw=modelblue,rounded corners,align=center,text width=3cm] (answer) at (9.6,0)
    {Both nonnegative?\\Keep the basis\\and update the result};
    \draw[model flow] (data)--(tests);
    \draw[model flow] (tests)--(answer);
  \end{tikzpicture}
  \caption{An optimal basis is a reusable certificate. The updated basis matrix
    must first be nonsingular.}\label{fig:certificate}
\end{figure}

\section{Changing one right-hand side}\label{sec:rhs}
Change only equation $r$:
\[
  b'=b+t e_r,\qquad v=A_B^{-1}e_r,\qquad \bar b'=\bar b+t v.
\]
Since $A$ and $c$ stay fixed, $q$ and $y$ stay fixed. The basis remains optimal
exactly when every basic value stays nonnegative:
\begin{equation}\label{eq:rhs-interval}
  \max_{i:v_i>0}\left(-\frac{\bar b_i}{v_i}\right)
  \leq t\leq
  \min_{i:v_i<0}\left(-\frac{\bar b_i}{v_i}\right).
\end{equation}
An empty maximum is $-\infty$, and an empty minimum is $+\infty$.
A zero $v_i$ imposes no new restriction because the starting $\bar b_i\geq0$.
Indices $i$ refer to \emph{basis rows}; $r$ refers to an original equation.

Within this closed interval,
\begin{equation}\label{eq:shadow}
  x_B(t)=\bar b+t v,\qquad
  z^*(b+t e_r)=z_0+t c_B^Tv=z_0+t y_r.
\end{equation}
The multiplier $y_r$ is the \emph{shadow price}: objective change per unit of
RHS change, while this basis is valid. At a finite endpoint a basic variable
becomes zero. The same basis still certifies optimality there, but degeneracy
and a different marginal value beyond the endpoint are possible.

Compute $v$ by solving $A_Bv=e_r$. A final tableau column gives $v$ only if its
original column is $+e_r$. A surplus column $-e_r$ gives $-v$.
An artificial column is not part of the final original LP; retain the row
transformation explicitly rather than relying on a deleted artificial.

\subsection{A worked sensitivity tableau}\label{sec:main-example}
Consider
\begin{equation}\label{eq:main-lp}
  \begin{aligned}
    \max\quad&-x_1+3x_2-2x_3\\
    \text{subject to}\quad&3x_1-x_2+2x_3\leq7,\\
    &-2x_1+4x_2\leq12,\\
    &-4x_1+3x_2+8x_3\leq10,\qquad x_1,x_2,x_3\geq0.
  \end{aligned}
\end{equation}
Ratios are calculated in the column to the right of the RHS;
$\longrightarrow$ in the separate Out column marks the chosen leaving row.
The indicator row is labelled $z_j-c_j$; $q_j$ abbreviates this difference
in formulas. Vertical rules separate $c_B$, the basic-variable labels,
the coefficient columns, and the RHS. Variable-column headings keep their
order across pivots; only the basic-variable labels change. Artificial
columns are removed after Phase I without reordering the retained columns.

Introduce slacks $x_4,x_5,x_6$. For $B=(1,2,6)$ the final tableau is
\begin{center}
  \begin{tabular}{r@{\hspace{0.2em}}|@{\hspace{0.2em}}c@{\enspace}|@{\enspace}rrrrrr@{\quad}|@{\quad}r}\toprule
    $c_B$ & Basic & $x_1$ & $x_2$ & $x_3$ & $x_4$ & $x_5$ & $x_6$ & RHS\\
    & $c_j$ & $-1$ & $3$ & $-2$ & $0$ & $0$ & $0$ & \\
    \midrule
    $-1$ & $x_1$ & 1 & 0 & $4/5$ & $2/5$ & $1/10$ & 0 & 4\\
    $3$ & $x_2$ & 0 & 1 & $2/5$ & $1/5$ & $3/10$ & 0 & 5\\
    $0$ & $x_6$ & 0 & 0 & 10 & 1 & $-1/2$ & 1 & 11\\
    \midrule
    & $z_j$ & $-1$ & $3$ & $\tfrac{2}{5}$ & $\tfrac{1}{5}$ & $\tfrac{4}{5}$ & $0$ & $11$\\
    & $z_j-c_j$ & 0 & 0 & $12/5$ & $1/5$ & $4/5$ & 0 & $z_0=11$\\
    \bottomrule
  \end{tabular}
\end{center}
This gives $x^*=(4,5,0)$, slacks $(0,0,11)$, and
$y=(1/5,4/5,0)$. The basis is nondegenerate and all nonbasic indicators are
strictly positive, so the optimizer is unique.

Change the second RHS to $12+t$. The $x_5$ column is $v$, hence
\[
  x_B(t)=\begin{pmatrix}4+t/10\\5+3t/10\\11-t/2\end{pmatrix},
  \qquad \max\{-40,-50/3\}\leq t\leq22.
\]
Therefore
\begin{equation}\label{eq:rhs-result}
  -\frac{50}{3}\leq t\leq22,\qquad
  -\frac{14}{3}\leq b_2'\leq34,\qquad z^*=11+\frac45t.
\end{equation}
The negative lower bound on $b_2'$ is allowed: the row has a negative $x_1$
coefficient, so a negative RHS need not make it infeasible.

\begin{figure}[H]\centering
  \begin{tikzpicture}[x=0.25cm,y=0.21cm,font=\small]
    \fill[modelblue!7] (-16.6667,0) rectangle (22,21);
    \draw[model axis] (-20,0)--(27,0) node[right] {$t$};
    \draw[model axis] (0,-1)--(0,23) node[above] {basic value};
    \draw[modelblue,thick] (-16.6667,2.3333)--(22,6.2) node[right] {$x_1$};
    \draw[modelorange,thick] (-16.6667,0)--(22,11.6) node[right] {$x_2$};
    \draw[black,thick] (-16.6667,19.3333)--(22,0);
    \node at (-8,16.8) {$x_6$};
    \draw[dashed] (-16.6667,0)--(-16.6667,21);
    \draw[dashed] (22,0)--(22,21);
    \node[below] at (-16.6667,0) {$-50/3$};
    \node[below] at (22,0) {$22$};
    \node[below left] at (0,0) {$0$};
    \node[above] at (3,21) {all three values nonnegative};
  \end{tikzpicture}
  \caption{The allowable interval is the intersection of three basic-value
    inequalities. Its endpoints are determined by $x_2=0$ and $x_6=0$.}
\end{figure}

\subsection{Endpoints, exercises, and the whole value function}
\begin{center}
  \begin{tabular}{r|rrr|r}\toprule
    $b_2'$ & $x_1$ & $x_2$ & $x_6$ & $z^*$\\\midrule
    $-14/3$ & $7/3$ & 0 & $58/3$ & $-7/3$\\
    12 & 4 & 5 & 11 & 11\\
    34 & $31/5$ & $58/5$ & 0 & $143/5$\\\bottomrule
  \end{tabular}
\end{center}
The endpoint calculation answers the exercise:
\begin{quote}\selectlanguage{polish}
  Jakie jest RO zadania, gdy zmienimy $b_2=-14/3$?
\end{quote}
\selectlanguage{english}
Here RO means an optimal solution. Substitute $t=-50/3$ in the basic vector,
rather than substituting $-14/3$ for the \emph{change}.

For the first RHS the $x_4$ column gives
\[
  (x_1,x_2,x_6)=(4+2t/5,\ 5+t/5,\ 11+t),\qquad
  t\geq-10,\quad b_1'\geq-3,\quad z^*=11+t/5.
\]
Thus the answer to \emph{Optymalna wartość f. celu ulega zmianie. O ile?}
is $\Delta z=t/5$. The third RHS can change by $t\geq-11$:
only $x_6=11+t$ changes, and the value remains 11.

We can go beyond the basis interval for the second RHS, denoted now by $r$:
\begin{equation}\label{eq:value-function}
  z^*(r)=
  \begin{cases}
    \text{infeasible}, & r<-14/3,           \\
    (4r+7)/5,          & -14/3\leq r\leq34, \\
    143/5,             & r\geq34.
  \end{cases}
\end{equation}
To justify all three cases, first observe that $x_3$ can be set to zero:
doing so relaxes constraints and increases the objective if $x_3>0$.
The first constraint gives $x_1\leq(7+x_2)/3$, so
$-2x_1+4x_2\geq-14/3+(10/3)x_2\geq-14/3$.
Below $-14/3$ feasibility is impossible. In the middle range,
\eqref{eq:rhs-result} supplies a feasible optimum and a dual certificate.
For $r\geq34$, $x=(31/5,58/5,0)$ is feasible, and
$y=(9/5,0,8/5)$ is dual feasible with value $143/5$.
Matching values prove the final branch.

\begin{figure}[H]\centering
  \begin{tikzpicture}[x=0.23cm,y=0.14cm,font=\small]
    \fill[pattern=north east lines,pattern color=gray!60] (-10,-5) rectangle (-4.6667,32);
    \draw[model axis] (-11,0)--(49,0) node[right] {$r=b_2'$};
    \draw[model axis] (0,-5)--(0,34) node[above] {$z^*(r)$};
    \draw[modelblue,very thick] (-4.6667,-2.3333)--(34,28.6)--(46,28.6);
    \draw[dashed] (34,0)--(34,28.6);
    \draw[dashed] (12,0)--(12,11);
    \node[model point] at (-4.6667,-2.3333) {};
    \node[model point] at (12,11) {};
    \node[model point] at (34,28.6) {};
    \node[below] at (-4.6667,-5) {$-14/3$};
    \node[below] at (12,0) {$12$};
    \node[below] at (34,0) {$34$};
    \node[above,rotate=26] at (19,20) {slope $4/5$};
    \node[above] at (40,28.6) {slope $0$};
    \node[rotate=90,fill=white,inner sep=2pt] at (-7,17) {infeasible};
  \end{tikzpicture}
  \caption{A shadow price applies to a branch of the value function.
    After $r=34$, the second resource has zero marginal value.}
\end{figure}

\section{Changing one objective coefficient}\label{sec:cost}
Changing $c$ preserves every feasible point and the old basic vector.
It changes the dual certificate and, usually, the value.

\subsection{A nonbasic coefficient}
If $k\in N$ and $c_k'=c_k+t$, then $c_B$ and $y$ are unchanged:
\[
  q_k'=q_k-t,\qquad q_j'=q_j\ (j\ne k).
\]
The basis remains optimal for $t\leq q_k$. Its objective is unchanged because
$x_k=0$. For \eqref{eq:main-lp},
\[
  t\leq12/5\quad\Longleftrightarrow\quad c_3'\leq2/5.
\]
This answers \emph{O ile może wzrosnąć wsp. $c_3$, nie wpływając na zmianę RO?}
The maximum increase is $12/5$, not the new coefficient $2/5$.

At equality a zero indicator appears. To establish \emph{distinct} optimizers,
also check that a positive feasible step is possible. Here $x_3=\theta$ gives
\[
  (x_1,x_2,x_6)=(4-4\theta/5,\ 5-2\theta/5,\ 11-10\theta),
  \qquad 0\leq\theta\leq11/10.
\]
For $c_3'=2/5$ this whole segment has value 11. For a larger $c_3'$,
even a small positive step improves the objective, so the old point is no
longer optimal.

\subsection{A basic coefficient}
If $k=j_s\in B$ and $c_k'=c_k+t$, then $c_B'=c_B+t e_s$ and
\[
  q_j'=q_j+t\bar a_{sj}\qquad(j\in N).
\]
For basic columns the simultaneous change in $c_k$ cancels, so $q_B'=0$.
The exact basis range is
\begin{equation}\label{eq:basic-cost-range}
  \max_{j\in N:\bar a_{sj}>0}\left(-\frac{q_j}{\bar a_{sj}}\right)
  \leq t\leq
  \min_{j\in N:\bar a_{sj}<0}\left(-\frac{q_j}{\bar a_{sj}}\right).
\end{equation}
Use original nonbasic columns only, not basic or removed artificial columns.
A zero coefficient adds no restriction when $q_j\geq0$.
The old point remains fixed and $z'=z_0+t x_k$.

For the main example:
\begin{center}
  \begin{tabular}{c|rrr|c|c}\toprule
    Changed cost & From $x_3$ & From $x_4$ & From $x_5$
    & Allowable change & Value\\\midrule
    $c_1=-1+t$ & $t\geq-3$ & $t\geq-1/2$ & $t\geq-8$
    & $t\geq-1/2$ & $11+4t$\\
    $c_2=3+t$ & $t\geq-6$ & $t\geq-1$ & $t\geq-8/3$
    & $t\geq-1$ & $11+5t$\\\bottomrule
  \end{tabular}
\end{center}
Thus $c_1'\geq-3/2$ and $c_2'\geq2$, with no upper bound.
At $t=-1/2$ or $-1$ a nonbasic indicator vanishes and a positive step is
available; alternative optimal points appear. Away from these endpoints,
the original strict indicators give uniqueness.

\begin{figure}[H]\centering
  \begin{tikzpicture}[x=2.1cm,y=0.62cm,font=\small]
    \fill[modelblue!7] (-0.5,0) rectangle (2.5,5.6);
    \draw[model axis] (-1.5,0)--(3,0) node[right] {$t=\Delta c_1$};
    \draw[model axis] (0,-0.5)--(0,6) node[above] {$q_j(t)$};
    \draw[modelblue,thick] (-1,1.6)--(2.5,4.4) node[right] {$q_3$};
    \draw[modelorange,thick] (-0.5,0)--(2.5,1.2) node[right] {$q_4$};
    \draw[black,thick] (-1,0.7)--(2.5,1.05);
    \node[above] at (-0.95,0.7) {$q_5$};
    \draw[dashed] (-0.5,0)--(-0.5,5.6);
    \node[below] at (-0.5,0) {$-1/2$};
    \node[above] at (1,5.6) {all indicators nonnegative};
  \end{tikzpicture}
  \caption{For a basic-cost change, several indicators move together.
    Here $q_4(t)=1/5+2t/5$ determines the endpoint.}
\end{figure}

\section{Vector changes and parameter regions}\label{sec:vector}
With fixed $A$, simultaneous changes $b'=b+\delta b$ and $c'=c+\delta c$
give
\begin{equation}\label{eq:joint-tests}
  \bar b'=\bar b+A_B^{-1}\delta b,\qquad
  q'=q+\bar A^T\delta c_B-\delta c.
\end{equation}
These are direct, coupled tests of the old basis. When both pass,
\begin{equation}\label{eq:joint-value}
  z'=(c_B+\delta c_B)^TA_B^{-1}(b+\delta b).
\end{equation}
If only $b$ changes, this reduces to $z_0+y^T\delta b$.
If only $c$ changes, it reduces to $z_0+x^T\delta c$.
If both change, include the cross term
$\delta c_B^TA_B^{-1}\delta b$.

One-at-a-time intervals cannot be combined independently. In
\eqref{eq:main-lp}, $\delta b=(0,20,0)$ gives $x_B=(6,11,1)$ and
$\delta b=(0,0,-10)$ gives $(4,5,1)$: both are allowable separately.
Together they give $(6,11,-9)$, which is infeasible.
The exact joint region, rather than a rectangular collection of intervals,
settles the question.

\subsection{A smaller model for all subsequent updates}\label{sec:small-model}
We will repeatedly reuse
\begin{equation}\label{eq:small-lp}
  \max\ 2x_1-x_2+x_3,\qquad
  x_1+x_2+x_3\leq6,\quad -x_1+2x_2\leq4,\quad x\geq0.
\end{equation}
Slacks are $x_4,x_5$, $B=(1,5)$, and
\[
  A_B=\begin{pmatrix}1&0\\-1&1\end{pmatrix},\qquad
  A_B^{-1}=\begin{pmatrix}1&0\\1&1\end{pmatrix}.
\]
The inverse is displayed to understand the row transformation; numerical code
should solve linear systems.
\begin{center}
  \begin{tabular}{r@{\hspace{0.2em}}|@{\hspace{0.2em}}c@{\enspace}|@{\enspace}rrrrr@{\quad}|@{\quad}r}\toprule
    $c_B$ & Basic & $x_1$ & $x_2$ & $x_3$ & $x_4$ & $x_5$ & RHS\\
    & $c_j$ & $2$ & $-1$ & $1$ & $0$ & $0$ & \\
    \midrule
    $2$ & $x_1$ & 1 & 1 & 1 & 1 & 0 & 6\\
    $0$ & $x_5$ & 0 & 3 & 1 & 1 & 1 & 10\\
    \midrule
    & $z_j$ & $2$ & $2$ & $2$ & $2$ & $0$ & $12$\\
    & $z_j-c_j$ & 0 & 3 & 1 & 2 & 0 & $z_0=12$\\
    \bottomrule
  \end{tabular}
\end{center}
The unique decision optimum is $(6,0,0)$; $y=(2,0)$.

For arbitrary new RHS $(u,v)$ the same basis is optimal exactly when
$u\geq0$ and $u+v\geq0$. The basic values are $(u,u+v)$ and the value is $2u$.
For arbitrary decision costs $(\gamma_1,\gamma_2,\gamma_3)$ with slack costs
zero, the basis is optimal exactly when
\[
  q_N=(\gamma_1-\gamma_2,\ \gamma_1-\gamma_3,\ \gamma_1)\geq0.
\]
This means $\gamma_1\geq0$, $\gamma_2\leq\gamma_1$, and
$\gamma_3\leq\gamma_1$. These regions live in \emph{data space}, not
decision space.

\begin{figure}[H]\centering
  \begin{tikzpicture}[font=\small]
    \begin{scope}[x=0.65cm,y=0.45cm]
      \fill[modelblue!12] (0,0)--(7,-7)--(7,6)--(0,6)--cycle;
      \draw[modelblue,thick] (0,0)--(7,-7);
      \draw[modelblue,thick] (0,0)--(0,6);
      \draw[model axis] (-0.6,0)--(7.6,0) node[right] {$u$};
      \draw[model axis] (0,-7.5)--(0,6.7) node[above] {$v$};
      \node[model point] at (6,4) {};
      \node[above left] at (6,4) {$(6,4)$};
      \node[model point] at (3,4) {};
      \node[above left] at (3,4) {$(3,4)$};
      \node[rotate=-35,below] at (4,-4) {$u+v=0$};
      \node[align=center] at (4,-0.9) {same basis\\$u\geq0,\ u+v\geq0$};
    \end{scope}
    \begin{scope}[shift={(8.1,0)},x=0.75cm,y=0.55cm]
      \fill[modelblue!12] (1,-4)--(5,-4)--(5,5)--(1,1)--cycle;
      \draw[modelblue,thick] (1,-4)--(1,1)--(5,5);
      \draw[model axis] (-0.5,0)--(5.7,0) node[right] {$\gamma_1$};
      \draw[model axis] (0,-4.7)--(0,5.7) node[above] {$\gamma_2$};
      \node[model point] at (2,-1) {};
      \node[below right] at (2,-1) {$(2,-1)$};
      \node[fill=modelorange,circle,inner sep=1.7pt] at (3,4) {};
      \node[above left] at (3,4) {$(3,4)$};
      \node[rotate=36,below] at (3.7,3.7) {$\gamma_2=\gamma_1$};
      \node[align=center] at (3,-3) {same basis\\slice $\gamma_3=1$};
    \end{scope}
  \end{tikzpicture}
  \caption{RHS feasibility and cost optimality regions for $B=(1,5)$.
    The shaded regions continue beyond the plotting windows.
    The orange cost point is outside; the two marked RHS points are inside.}
\end{figure}

\subsection{Reoptimize after changing the whole objective}
Set decision costs to $(3,4,2)$. The old basic point is still feasible, its value
is now 18, and its new indicator row is $(0,-1,1,3,0)$.
\begin{center}
  \begin{tabular}{r@{\hspace{0.2em}}|@{\hspace{0.2em}}c@{\enspace}|@{\enspace}rrrrr@{\quad}|@{\quad}r@{\qquad\quad}l@{\quad}c}\cmidrule(r{1.5em}){1-8}
    $c_B$ & Basic & $x_1$ & $x_2$ & $x_3$ & $x_4$ & $x_5$ & RHS & Ratio ($x_2\uparrow$) & Out\\
    & $c_j$ & $3$ & $4$ & $2$ & $0$ & $0$ & & & \\
    \cmidrule(r{1.5em}){1-8}
    $3$ & $x_1$ & 1 & 1 & 1 & 1 & 0 & 6 & $6/1=6$ & \\
    $0$ & $x_5$ & 0 & 3 & 1 & 1 & 1 & 10 & $10/3$ & $\longrightarrow$\\
    \cmidrule(r{1.5em}){1-8}
    & $z_j$ & $3$ & $3$ & $3$ & $3$ & $0$ & $18$ & & \\
    & $z_j-c_j$ & 0 & $-1$ & 1 & 3 & 0 & $z_0=18$ & & \\
    \cmidrule(r{1.5em}){1-8}
  \end{tabular}
\end{center}
Enter $x_2$ and remove $x_5$ with step $10/3$.
\begin{center}
  \begin{tabular}{r@{\hspace{0.2em}}|@{\hspace{0.2em}}c@{\enspace}|@{\enspace}rrrrr@{\quad}|@{\quad}r}\toprule
    $c_B$ & Basic & $x_1$ & $x_2$ & $x_3$ & $x_4$ & $x_5$ & RHS\\
    & $c_j$ & $3$ & $4$ & $2$ & $0$ & $0$ & \\
    \midrule
    $3$ & $x_1$ & 1 & 0 & $2/3$ & $2/3$ & $-1/3$ & $8/3$\\
    $4$ & $x_2$ & 0 & 1 & $1/3$ & $1/3$ & $1/3$ & $10/3$\\
    \midrule
    & $z_j$ & $3$ & $4$ & $\tfrac{10}{3}$ & $\tfrac{10}{3}$ & $\tfrac{1}{3}$ & $\tfrac{64}{3}$\\
    & $z_j-c_j$ & 0 & 0 & $4/3$ & $10/3$ & $1/3$ & $z_0=64/3$\\
    \bottomrule
  \end{tabular}
\end{center}
The new optimum is $(8/3,10/3,0)$, with dual $y=(10/3,1/3)$.
The gain from the \emph{new starting value} is $10/3$, not $64/3-12$.

\begin{figure}[H]\centering
  \begin{tikzpicture}[x=0.85cm,y=0.85cm,font=\small]
    \draw[model region] (0,0)--(6,0)--(2.6667,3.3333)--(0,2)--cycle;
    \draw[model axis] (-0.3,0)--(6.8,0) node[right] {$x_1$};
    \draw[model axis] (0,-0.3)--(0,4.6) node[above] {$x_2$};
    \draw[model objective] (3.5556,2.6667)--(1.7778,4);
    \node[model point] at (6,0) {};
    \node[below] at (6,0) {old $(6,0)$};
    \node[model point] at (2.6667,3.3333) {};
    \node[above right] at (2.6667,3.3333) {new $(8/3,10/3)$};
    \draw[model flow] (5.5,0.5)--(3.2,2.8);
    \node[align=center] at (7.3,2) {$x_3=0$ slice\\new objective\\$3x_1+4x_2$};
  \end{tikzpicture}
  \caption{Changing the objective leaves the feasible set fixed and changes the
    supporting line. A single primal pivot follows the improving edge.}
\end{figure}

\subsection{Change the RHS vector}
Set $b'=(3,4)$. The new basic vector is $(3,7)$ and $q$ is unchanged:
$x=(3,0,0)$ is optimal, with value 6. No pivot is needed.

A failed basis test can also expose genuine infeasibility. For $b'=(6,-8)$,
the basic values are $(6,-2)$ and the second canonical equation is
\[
  3x_2+x_3+x_4+x_5=-2.
\]
Its left side is nonnegative. The contradiction certifies infeasibility;
there is no eligible negative entry for a dual-simplex pivot.

\section{Choosing a method after a model change}\label{sec:methods}
\begin{center}
  \begin{tabular}{>{\raggedright\arraybackslash}p{3.4cm}>{\raggedright\arraybackslash}p{4.6cm}>{\raggedright\arraybackslash}p{5.1cm}}
    \toprule
    Change & What the old basis keeps & What to do if its test fails\\\midrule
    Costs $c$ & Primal feasibility & Recompute $q$; primal simplex\\
    RHS $b$ & Dual feasibility & Recompute $\bar b$; dual simplex\\
    New nonnegative variable & Primal feasibility at new variable zero & Test its indicator; primal simplex\\
    New $\leq$ constraint & Dual feasibility with a basic slack & Canonicalize the row; dual simplex\\
    Both $b$ and $c$ & Neither guaranteed & Test both; choose a valid start\\
    Matrix $A$ & Neither guaranteed & Check nonsingularity, then both tests\\\bottomrule
  \end{tabular}
\end{center}
This is a choice of starting method, not a guarantee that the modified problem
has a finite optimum. Primal simplex can expose an improving ray;
dual simplex can expose a contradictory row.

For a changed matrix $A'$, the old column indices remain usable only if $A_B'$
is nonsingular. Even one changed basic column can affect every transformed
coefficient. If only a nonbasic column changes, $\bar b$ is unchanged and
only that column's indicator needs retesting (when costs are unchanged).
Use fresh linear solves rather than extending a one-parameter formula beyond
its hypotheses.

\section{Adding a nonnegative decision variable}\label{sec:new-variable}
Add $w\geq0$ with column $a$ and objective coefficient $\gamma$:
\[
  Ax+aw=b,\qquad \max(c^Tx+\gamma w).
\]
At $w=0$ the old basic point is feasible. Append
\begin{equation}\label{eq:new-column}
  u=A_B^{-1}a,\qquad q_w=c_B^Tu-\gamma=y^Ta-\gamma.
\end{equation}
If $q_w\geq0$, the old basis certifies optimality with $w=0$.
If $q_w<0$, it loses its dual certificate; primal simplex can enter $w$.
From a nondegenerate feasible basis, a positive feasible step improves the
objective. For a degenerate basis the first step can be zero; reoptimization
determines whether the old point is still optimal.
If $u\leq0$, the feasible direction $d_w=1$, $d_B=-u$, $d_N=0$
gives an unbounded improving ray.

For resource inequalities, $y^Ta$ is the resource value of the new activity.
Compare it with the activity's revenue $\gamma$. A negative component of $a$
means the activity releases that resource, rather than consuming it.

\begin{figure}[H]\centering
  \begin{tikzpicture}[font=\small]
    \node[draw=modelblue,align=center,text width=3.1cm] (a) at (0,0)
    {New activity\\column $a$, return $\gamma$};
    \node[draw=modelblue,align=center,text width=3.6cm] (price) at (4.8,0)
    {Current resource prices $y$\\value $y^Ta$};
    \node[draw=modelorange,align=center,text width=3.5cm] (gain) at (9.8,0)
    {Net marginal gain\\$\gamma-y^Ta=-q_w$};
    \draw[model flow] (a)--(price);
    \draw[model flow] (price)--(gain);
  \end{tikzpicture}
  \caption{Pricing a new column gives an immediate test of the old dual
    certificate. Prices describe the current basis, not all later pivots.}
\end{figure}

\subsection{Worked column update}
For \eqref{eq:small-lp}, introduce $x_6\geq0$ with
\[
  a=\begin{pmatrix}-1\\2\end{pmatrix},\qquad \gamma=1,\qquad
  u=\begin{pmatrix}-1\\1\end{pmatrix},\qquad q_6=-3.
\]
The canonical equations are now
\[
  x_1+x_2+x_3+x_4-x_6=6,\qquad
  3x_2+x_3+x_4+x_5+x_6=10.
\]
Enter $x_6$. Only the second row has a positive coefficient; the step is 10
and $x_5$ leaves. The final tableau is
\begin{center}
  \begin{tabular}{r@{\hspace{0.2em}}|@{\hspace{0.2em}}c@{\enspace}|@{\enspace}rrrrrr@{\quad}|@{\quad}r}\toprule
    $c_B$ & Basic & $x_1$ & $x_2$ & $x_3$ & $x_4$ & $x_5$ & $x_6$ & RHS\\
    & $c_j$ & $2$ & $-1$ & $1$ & $0$ & $0$ & $1$ & \\
    \midrule
    $2$ & $x_1$ & 1 & 4 & 2 & 2 & 1 & 0 & 16\\
    $1$ & $x_6$ & 0 & 3 & 1 & 1 & 1 & 1 & 10\\
    \midrule
    & $z_j$ & $2$ & $11$ & $5$ & $5$ & $3$ & $1$ & $42$\\
    & $z_j-c_j$ & 0 & 12 & 4 & 5 & 3 & 0 & $z_0=42$\\
    \bottomrule
  \end{tabular}
\end{center}
The extended optimum is $(16,0,0,0,0,10)$, with objective $2(16)+10=42$.
The gain is $(-q_6)10=30$.
Both original inequalities with the new column are tight, and
$y=(5,3)$ is dual feasible with value $6(5)+4(3)=42$.

\section{Adding a constraint}\label{sec:new-constraint}
A new constraint shrinks the feasible set. If an old optimizer satisfies it,
that point remains optimal: no point in the smaller set can beat the old
maximum. If the old point violates it, a new feasible point is required.

For a new inequality $a^Tx\leq\beta$, introduce a slack $s\geq0$:
$a^Tx+s=\beta$. Substitute $x_B=\bar b-\bar A_Nx_N$ to obtain
\begin{equation}\label{eq:new-row}
  s+(a_N^T-a_B^T\bar A_N)x_N=\beta-a_B^T\bar b.
\end{equation}
The old basic columns have been eliminated from this row. With $s$ basic and
zero cost, the enlarged basis is nonsingular and its indicator row is
$(q,0)$. The new slack value is the constraint residual at the old point.
A negative residual gives a dual feasible, primal infeasible start for
dual simplex. Merely appending the untransformed row does not produce a
canonical tableau.

For a new equality, no permanent slack can absorb a residual. If the old point
violates it, one option is a signed artificial:
$a^Tx+\sigma w=\beta$, $w\geq0$, where $\sigma$ is chosen so that
$w=|\beta-a^Tx^*|$ at the old point. Phase I minimizes $w$.
A positive Phase I minimum certifies infeasibility. At zero, remove the
artificial and restore the true objective; a zero basic artificial needs
pivoting out or redundant-row cleanup. Equalities already satisfied by the
old point need no change of optimizer, though a basis representation may
need adjustment.

\subsection{Worked row update and a dual pivot}
Add $x_1-2x_3\leq-2$ to \eqref{eq:small-lp}.
The old point $(6,0,0)$ violates the row. With new slack $x_6$,
\[
  x_1-2x_3+x_6=-2.
\]
Subtract the old $x_1$ row. The new canonical tableau is
\begin{center}
  \begin{tabular}{r@{\hspace{0.2em}}|@{\hspace{0.2em}}c@{\enspace}|@{\enspace}rrrrrr@{\quad}|@{\quad}r@{\qquad\quad}l@{\quad}c}\cmidrule(r{1.5em}){1-9}
    $c_B$ & Basic & $x_1$ & $x_2$ & $x_3$ & $x_4$ & $x_5$ & $x_6$ & RHS & Dual ratios ($r=3$) & Out\\
    & $c_j$ & $2$ & $-1$ & $1$ & $0$ & $0$ & $0$ & & & \\
    \cmidrule(r{1.5em}){1-9}
    $2$ & $x_1$ & 1 & 1 & 1 & 1 & 0 & 0 & 6 & $x_2:\ 3/(1)=3$ & \\
    $0$ & $x_5$ & 0 & 3 & 1 & 1 & 1 & 0 & 10 & $x_3:\ 1/(3)=1 / 3$ & \\
    $0$ & $x_6$ & 0 & $-1$ & $-3$ & $-1$ & 0 & 1 & $-8$ & $x_4:\ 2/(1)=2$ & $\longrightarrow$\\
    \cmidrule(r{1.5em}){1-9}
    & $z_j$ & $2$ & $2$ & $2$ & $2$ & $0$ & $0$ & $12$ & & \\
    & $z_j-c_j$ & 0 & 3 & 1 & 2 & 0 & 0 & $z_0=12$ & & \\
    \cmidrule(r{1.5em}){1-9}
  \end{tabular}
\end{center}
Select the negative RHS row and compare the dual ratios displayed to the right.
Enter $x_3$, leave $x_6$, and pivot on $-3$:
\begin{center}
  \begin{tabular}{r@{\hspace{0.2em}}|@{\hspace{0.2em}}c@{\enspace}|@{\enspace}rrrrrr@{\quad}|@{\quad}r}\toprule
    $c_B$ & Basic & $x_1$ & $x_2$ & $x_3$ & $x_4$ & $x_5$ & $x_6$ & RHS\\
    & $c_j$ & $2$ & $-1$ & $1$ & $0$ & $0$ & $0$ & \\
    \midrule
    $2$ & $x_1$ & 1 & $2/3$ & 0 & $2/3$ & 0 & $1/3$ & $10/3$\\
    $0$ & $x_5$ & 0 & $8/3$ & 0 & $2/3$ & 1 & $1/3$ & $22/3$\\
    $1$ & $x_3$ & 0 & $1/3$ & 1 & $1/3$ & 0 & $-1/3$ & $8/3$\\
    \midrule
    & $z_j$ & $2$ & $\tfrac{5}{3}$ & $1$ & $\tfrac{5}{3}$ & $0$ & $\tfrac{1}{3}$ & $\tfrac{28}{3}$\\
    & $z_j-c_j$ & 0 & $8/3$ & 0 & $5/3$ & 0 & $1/3$ & $z_0=28/3$\\
    \bottomrule
  \end{tabular}
\end{center}
The optimum is $(10/3,0,8/3)$ with value $28/3$.
Its slacks are $(0,22/3,0)$; $y=(5/3,0,1/3)$ gives the matching
dual value $6(5/3)+4(0)-2(1/3)=28/3$.
The value decreases from 12, as a smaller feasible set requires.

\begin{figure}[H]\centering
  \begin{tikzpicture}[x=0.82cm,y=0.75cm,font=\small]
    \draw[fill=gray!8,draw=gray] (0,0)--(6,0)--(0,6)--cycle;
    \draw[model region] (0,1)--(0,6)--(3.3333,2.6667)--cycle;
    \draw[model axis] (-0.3,0)--(6.8,0) node[right] {$x_1$};
    \draw[model axis] (0,-0.3)--(0,6.7) node[above] {$x_3$};
    \draw[modelorange,thick] (0,1)--(5.5,3.75);
    \node[above,rotate=25] at (4.7,3.35) {$x_1-2x_3=-2$};
    \node[model point] at (6,0) {};
    \node[below] at (6,0) {old point};
    \node[model point] at (3.3333,2.6667) {};
    \node[below left] at (3.3333,2.6667) {new point};
    \node at (0.8,3) {$\Omega'$};
    \node[align=center] at (6.8,5.5) {$x_2=0$ slice\\gray: old set\\blue: new set};
  \end{tikzpicture}
  \caption{The new inequality cuts off the old optimizer.
    This is the $x_2=0$ slice; the tableau certifies the optimum in the full space.}
\end{figure}

\section{Removing a variable}\label{sec:delete-variable}
Removing $x_s$ means imposing $x_s=0$ and deleting its column and cost.
The feasible set, embedded in the old space, becomes smaller.
If an old optimizer has $x_s=0$, deleting that coordinate gives an optimizer
of the reduced problem. A nonbasic zero column can simply be removed.
A \emph{basic} zero variable requires a basis exchange or redundant-row
cleanup before its column is removed.

If $x_s^*>0$, feasibility of the reduced model is not assured.
A reliable feasibility phase is
\[
  \min x_s,\qquad Ax=b,\quad x\geq0.
\]
The old point provides a feasible start. A strictly positive minimum proves
that setting $x_s=0$ is impossible. At a zero minimum, remove the column after
cleaning up the basis, then restore and optimize the reduced objective.
A zero-valued variable can still be basic: basis membership alone is not an
infeasibility certificate. An unqualified finite Big M is not a substitute
for this feasibility test.

Because the original maximum is finite, any feasible reduced model is bounded
above by it. For a nonempty polyhedron, a bounded linear objective attains its
maximum. Here the possible outcomes are therefore finite optimum or
infeasibility, rather than a newly unbounded maximum.

\begin{center}
  \begin{tabular}{>{\raggedright\arraybackslash}p{4cm}>{\raggedright\arraybackslash}p{4.3cm}>{\raggedright\arraybackslash}p{4.8cm}}
    \toprule
    Deletion & Feasible-set effect & Value effect for maximization\\\midrule
    Variable & Impose its value zero & Cannot increase; may become infeasible\\
    Constraint & Relax the restrictions & Cannot decrease; may become unbounded\\\bottomrule
  \end{tabular}
\end{center}
For \eqref{eq:small-lp}, deleting $x_2$ retains $(6,0,0)$ and value 12.
Deleting $x_1$ gives
\[
  \max(-x_2+x_3),\qquad x_2+x_3\leq6,\quad 2x_2\leq4,\quad x_2,x_3\geq0.
\]
Since $-x_2+x_3\leq x_2+x_3\leq6$, $(x_2,x_3)=(0,6)$ is optimal.
By contrast, in $x_1+x_2=1$, $x_2=0$, $x\geq0$, deletion of $x_1$
leaves the impossible equation $0=1$.

\section{Removing a constraint}\label{sec:delete-constraint}
Removing a constraint enlarges the feasible set. The old point stays feasible,
but may cease to be optimal; a finite old optimum can become unbounded.

\subsection{When the old dual certificate survives}
For $\max c^Tx$ subject to $Gx\leq b$, $x\geq0$, a dual certificate is
$y\geq0$, $G^Ty\geq c$. If some optimal dual vector has $y_r=0$, drop its
$r$th component. The remaining vector is dual feasible for the model without
row $r$, and its objective is unchanged. The old primal point still has the
same value, so it remains optimal.

Strict slack at an old optimal point implies $y_r=0$ by complementarity
and is sufficient for this conclusion. It is not necessary: a tight row can
also have zero multiplier. For equalities the multipliers are free; a zero
multiplier similarly permits deletion of that row.
Checking one chosen dual vector is a sufficient test; failure of that test
alone does not determine the best new value.

\subsection{Relaxing a row without throwing away the tableau}
Deleting equation $r$ can be represented by
\[
  Ax+e_r u-e_r v=b,\qquad u,v\geq0,\quad c_u=c_v=0.
\]
The difference $u-v$ can absorb any residual of the deleted equation.
The other equations are unchanged. Original feasible vectors of the reduced
model correspond exactly to extended feasible vectors after choosing $u,v$.
The new indicators are $q_u=y_r$ and $q_v=-y_r$.
If either is negative, primal simplex can improve the old solution.

For a converted $\leq$ row, its existing slack already has column $+e_r$.
Add only the opposite column $-e_r$ with zero cost.
For a converted $\geq$ row, the existing surplus is $-e_r$, so add $+e_r$.
Once solved, discard the auxiliary pair and the deleted equation.
This construction can also expose unboundedness.

\subsection{Production model: two different deletions}
Consider the production model
\begin{equation}\label{eq:production}
  \max\,3x_1+2x_2,\quad
  2x_1+x_2\leq10,\quad 3x_1+3x_2\leq24,\quad
  \tfrac32x_1\leq6,\quad x\geq0.
\end{equation}
The optimum is $(2,6)$, value 18, with slacks $(x_3,x_4,x_5)=(0,0,3)$
and dual prices $(1,1/3,0)$.

Removing the third constraint retains the optimum: its slack is positive,
its dual multiplier is zero, and the first two rows still bound the objective
by $3x_1+2x_2=(2x_1+x_2)+(x_1+x_2)\leq10+8=18$.

Removing the first constraint changes the result. The remaining rows imply
$x_1\leq4$ and $x_1+x_2\leq8$, giving
\[
  3x_1+2x_2=x_1+2(x_1+x_2)\leq4+16=20.
\]
The point $(4,4)$ attains 20. This is an exact bound, not just a graphical guess.

\begin{figure}[H]\centering
  \begin{tikzpicture}[font=\small]
    \foreach \panel/\heading in {0/Original,1/Remove row 3,2/Remove row 1}{
        \begin{scope}[shift={({\panel*4.7},0)},x=0.44cm,y=0.43cm]
          \node[above] at (4,9.8) {\heading};
          \ifnum\panel=0
            \draw[model region] (0,0)--(4,0)--(4,2)--(2,6)--(0,8)--cycle;
          \fi
          \ifnum\panel=1
            \draw[model region] (0,0)--(5,0)--(2,6)--(0,8)--cycle;
          \fi
          \ifnum\panel=2
            \draw[model region] (0,0)--(4,0)--(4,4)--(0,8)--cycle;
          \fi
          \draw[model axis] (-0.4,0)--(8.6,0) node[right] {$x_1$};
          \draw[model axis] (0,-0.4)--(0,9.3) node[above] {$x_2$};
          \draw[gray] (0,8)--(8,0);
          \ifnum\panel=2
            \draw[model objective] (2.5,6.25)--(5,2.5);
            \node[model point] at (4,4) {};
            \node[above right] at (4,4) {$(4,4)$};
            \node[below] at (4,-0.6) {$z^*=20$};
          \else
            \draw[gray] (1,8)--(5,0);
            \draw[model objective] (1,7.5)--(4,3);
            \node[model point] at (2,6) {};
            \node[above right] at (2,6) {$(2,6)$};
            \node[below] at (4,-0.6) {$z^*=18$};
          \fi
        \end{scope}}
  \end{tikzpicture}
  \caption{Constraint-removal geometry: deleting a slack row preserves the
    optimizer; deleting a priced row permits improvement. All three sets enlarge
    or stay fixed relative to the relevant retained inequalities.}
\end{figure}

Here is the full update for deleting row 1 using $x_6$ with column $-e_1$.
Start from the ordered basis $B=(2,5,1)$:
\begin{center}
  \begin{tabular}{r@{\hspace{0.2em}}|@{\hspace{0.2em}}c@{\enspace}|@{\enspace}rrrrrr@{\quad}|@{\quad}r@{\qquad\quad}l@{\quad}c}\cmidrule(r{1.5em}){1-9}
    $c_B$ & Basic & $x_1$ & $x_2$ & $x_3$ & $x_4$ & $x_5$ & $x_6$ & RHS & Ratio ($x_6\uparrow$) & Out\\
    & $c_j$ & $3$ & $2$ & $0$ & $0$ & $0$ & $0$ & & & \\
    \cmidrule(r{1.5em}){1-9}
    $2$ & $x_2$ & 0 & 1 & $-1$ & $2/3$ & 0 & 1 & 6 & $6 / 1=6$ & \\
    $0$ & $x_5$ & 0 & 0 & $-3/2$ & $1/2$ & 1 & $3/2$ & 3 & $3 / (3 / 2)=2$ & $\longrightarrow$\\
    $3$ & $x_1$ & 1 & 0 & 1 & $-1/3$ & 0 & $-1$ & 2 & $\text{--}$ & \\
    \cmidrule(r{1.5em}){1-9}
    & $z_j$ & $3$ & $2$ & $1$ & $\tfrac{1}{3}$ & $0$ & $-1$ & $18$ & & \\
    & $z_j-c_j$ & 0 & 0 & 1 & $1/3$ & 0 & $-1$ & $z_0=18$ & & \\
    \cmidrule(r{1.5em}){1-9}
  \end{tabular}
\end{center}
Enter $x_6$ and remove $x_5$ with step 2.
\begin{center}
  \begin{tabular}{r@{\hspace{0.2em}}|@{\hspace{0.2em}}c@{\enspace}|@{\enspace}rrrrrr@{\quad}|@{\quad}r}\toprule
    $c_B$ & Basic & $x_1$ & $x_2$ & $x_3$ & $x_4$ & $x_5$ & $x_6$ & RHS\\
    & $c_j$ & $3$ & $2$ & $0$ & $0$ & $0$ & $0$ & \\
    \midrule
    $2$ & $x_2$ & 0 & 1 & 0 & $1/3$ & $-2/3$ & 0 & 4\\
    $0$ & $x_6$ & 0 & 0 & $-1$ & $1/3$ & $2/3$ & 1 & 2\\
    $3$ & $x_1$ & 1 & 0 & 0 & 0 & $2/3$ & 0 & 4\\
    \midrule
    & $z_j$ & $3$ & $2$ & $0$ & $\tfrac{2}{3}$ & $\tfrac{2}{3}$ & $0$ & $20$\\
    & $z_j-c_j$ & 0 & 0 & 0 & $2/3$ & $2/3$ & 0 & $z_0=20$\\
    \bottomrule
  \end{tabular}
\end{center}
The auxiliary $x_6=2$ is precisely the violation $2(4)+4-10=2$
of the deleted inequality. Discarding $x_3,x_6$ gives the reduced solution
$(x_1,x_2,x_4,x_5)=(4,4,0,0)$.
The remaining-row multipliers $(2/3,2/3)$ have value
$24(2/3)+6(2/3)=20$.

Deleting a row can also cause unboundedness: in $\max x_1$ subject to
$x_1\leq1$, $x_1\geq0$, removal of $x_1\leq1$ leaves the improving ray
$x_1=\theta$, $\theta\geq0$.

\section{Lagrange multipliers and active constraints}\label{sec:lagrange}
Write all inequalities in the same orientation:
\[
  \max c^Tx,\qquad g_i(x)=a_i^Tx-b_i\leq0\ (i\in I),\qquad
  h_j(x)=a_j^Tx-b_j=0\ (j\in E).
\]
Here $x\in\R^n$ is otherwise unrestricted. Include variable bounds explicitly
among the inequalities: $x_k\geq0$ means $g(x)=-x_k\leq0$.
An inequality is \emph{active} at $\hat x$ if $g_i(\hat x)=0$.
Denote these indices by $\mathcal A(\hat x)$.

\subsection{Feasible directions and outward normals}
For a feasible point $\hat x$, a direction $d$ allows sufficiently small
positive steps $\hat x+\alpha d$ exactly when
\[
  a_i^Td\leq0\quad(i\in\mathcal A(\hat x)),\qquad
  a_j^Td=0\quad(j\in E).
\]
Inactive rows impose no restriction on a sufficiently small step: each has
a strictly positive slack. Because there are finitely many rows, one
positive step bound can be chosen for all of them.
Optimality is equivalent to $c^Td\leq0$ for every such direction.

If the objective normal can be represented as
\begin{equation}\label{eq:normal-representation}
  c=\sum_{i\in I}\lambda_i a_i+\sum_{j\in E}\nu_j a_j,\qquad
  \lambda_i\geq0,\quad
  \lambda_i=0\text{ for }i\notin\mathcal A(\hat x),
\end{equation}
then its product with every feasible direction is nonpositive.
Equality normals can have coefficients of either sign; inequality normals
must point outward with nonnegative weights.

\begin{figure}[H]\centering
  \begin{tikzpicture}[font=\small]
    \begin{scope}[x=0.55cm,y=0.5cm]
      \draw[model region] (0,0)--(4,0)--(4,2)--(2,6)--(0,8)--cycle;
      \draw[model axis] (-0.3,0)--(5.5,0) node[right] {$x_1$};
      \draw[model axis] (0,-0.3)--(0,9.5) node[above] {$x_2$};
      \node[model point] at (2,6) {};
      \draw[model flow] (2,6)--(3,4) node[right] {$d_1$};
      \draw[model flow] (2,6)--(0.6,7.4);
      \node[above] at (1.1,7.6) {$d_2$};
      \node[above] at (2,10.6) {feasible edge directions};
    \end{scope}
    \begin{scope}[shift={(7,0)},x=0.9cm,y=0.9cm]
      \fill[modelblue!10] (0,0)--(5,2.5)--(5,5)--cycle;
      \draw[model axis] (-0.3,0)--(5.7,0);
      \draw[model axis] (0,-0.3)--(0,5.7);
      \draw[modelblue,thick,-{Stealth}] (0,0)--(2,1) node[below right] {$a_1=(2,1)$};
      \draw[modelblue,thick,-{Stealth}] (0,0)--(3,3) node[above left] {$a_2=(3,3)$};
      \draw[modelorange,very thick,-{Stealth}] (0,0)--(3,2) node[right] {$c=(3,2)$};
      \draw[dashed] (2,1)--(3,2)--(1,1);
      \node[above] at (2.5,5.2) {normal cone at $(2,6)$};
      \node[below] at (2.5,-0.5) {$c=a_1+\tfrac13 a_2$};
    \end{scope}
  \end{tikzpicture}
  \caption{For the production optimum, $c$ lies between the active outward
    normals. Its products with both feasible edge directions are negative.
    The right panel is a space of normals, not another feasible region.}
\end{figure}

\subsection{The Lagrangian and the four KKT conditions}
It is convenient to minimize the negative objective:
\begin{equation}\label{eq:lagrangian}
  L(x,\lambda,\nu)=-c^Tx+\sum_{i\in I}\lambda_i g_i(x)
  +\sum_{j\in E}\nu_j h_j(x).
\end{equation}
Then the Karush--Kuhn--Tucker (KKT) conditions are
\begin{center}
  \begin{tabular}{ll}\toprule
    Condition & Formula\\\midrule
    Primal feasibility & $g_i(\hat x)\leq0,\quad h_j(\hat x)=0$\\
    Multiplier signs & $\lambda_i\geq0,\quad \nu_j\in\R$\\
    Complementarity & $\lambda_i g_i(\hat x)=0$\\
    Stationarity & $-c+\sum_i\lambda_i a_i+\sum_j\nu_j a_j=0$\\\bottomrule
  \end{tabular}
\end{center}
For this linear program, these conditions are necessary and sufficient for
a feasible point to be optimal.

For sufficiency, compare any feasible $x$ with $\hat x$:
\[
  c^T(x-\hat x)=\sum_i\lambda_i a_i^T(x-\hat x)
  +\sum_j\nu_j a_j^T(x-\hat x)\leq0.
\]
Equalities contribute zero; positive inequality multipliers have active rows,
so their contributions are nonpositive.
For necessity, apply linear-programming strong duality to an optimal feasible
point, including its bound inequalities. An optimal dual vector exists.
The primal--dual gap is a sum of nonnegative multiplier--slack products;
zero gap gives complementarity, and the dual equations give stationarity.
No independence of active inequality normals is required.

For a nonlinear problem, replace $c$ and $a_i$ by the appropriate gradients.
For a differentiable concave maximization objective with convex inequalities
and affine equalities, KKT conditions still suffice globally.
Necessity generally requires a constraint qualification, such as Slater's
condition (with the usual domain assumptions); convexity alone is insufficient.
For example, $\max x$ subject to $x^2\leq0$ has the unique optimum $x=0$,
but stationarity would demand $-1+\lambda(2x)=0$, which is impossible there.

\subsection{A complete mixed-sign multiplier certificate}\label{sec:kkt-example}
Consider the question:
\begin{quote}\selectlanguage{polish}
  Czy $\hat x=[0,\ 17/4,\ 0,\ 23/4]^T$ jest RO?
\end{quote}
\selectlanguage{english}
The model is
\begin{equation}\label{eq:kkt-model}
  \begin{aligned}
    \max\quad&x_1+3x_2-2x_3+x_4\\
    \text{subject to}\quad&3x_1+2x_2+x_3+x_4\leq20,\\
    &2x_1+x_2+x_4=10,\\
    &5x_1-2x_2+x_3+2x_4\geq3,\qquad x\geq0.
  \end{aligned}
\end{equation}
Orient the third row as $3-5x_1+2x_2-x_3-2x_4\leq0$.
Assign $\lambda_1$ to row 1, $\nu$ to the equality, $\lambda_3$ to the
oriented row 3, and $\mu_k$ to the bound $-x_k\leq0$.
At $\hat x$ the row values are $57/4$, 10, and 3.
\begin{center}
  \begin{tabular}{c|c|c|c}\toprule
    Constraint & Nonnegative slack & Active? & Multiplier\\\midrule
    Row 1 ($\leq20$) & $23/4$ & No & $\lambda_1=0$\\
    Equality row & 0 & Yes & $\nu=2$ (free)\\
    Row 3 ($\geq3$) & 0 & Yes & $\lambda_3=1/2$\\
    $x_1\geq0$ & 0 & Yes & $\mu_1=1/2$\\
    $x_2\geq0$ & $17/4$ & No & $\mu_2=0$\\
    $x_3\geq0$ & 0 & Yes & $\mu_3=3/2$\\
    $x_4\geq0$ & $23/4$ & No & $\mu_4=0$\\\bottomrule
  \end{tabular}
\end{center}
Inactive-row multipliers vanish by complementarity.
Stationarity in coordinates 2 and 4 gives
\[
  -3+\nu+2\lambda_3=0,\qquad -1+\nu-2\lambda_3=0.
\]
Hence $\nu=2$ and $\lambda_3=1/2$. The remaining coordinates give
\[
  -1+2\nu-5\lambda_3-\mu_1=0,\qquad
  2-\lambda_3-\mu_3=0,
\]
so $\mu_1=1/2$ and $\mu_3=3/2$. All inequality multipliers are nonnegative.
Equivalently, the full normal identity is
\[
  \begin{pmatrix}1\\3\\-2\\1\end{pmatrix}
  =2\begin{pmatrix}2\\1\\0\\1\end{pmatrix}
  +\frac12\begin{pmatrix}-5\\2\\-1\\-2\end{pmatrix}
  -\frac12 e_1-\frac32 e_3.
\]
For any feasible $x$ this yields the explicit upper bound
\[
  c^Tx
  =2(2x_1+x_2+x_4)
  +\frac12(-5x_1+2x_2-x_3-2x_4)
  -\frac12x_1-\frac32x_3
  \leq20-\frac32=\frac{37}{2}.
\]
The candidate attains $37/2$, so it is optimal.
It is also unique: equality in this bound forces $x_1=x_3=0$ and the
third row tight. Together with $x_2+x_4=10$ these give
$x_2=17/4$, $x_4=23/4$.

\begin{figure}[H]\centering
  \begin{tikzpicture}[x=1.35cm,y=0.28cm,font=\small]
    \draw[model axis] (-0.2,0)--(5.8,0) node[right] {$x_2$};
    \draw[model axis] (0,-0.3)--(0,10.8) node[above] {$x_4$};
    \draw[gray,thick] (0,10)--(5,5);
    \draw[modelblue,very thick] (0,10)--(4.25,5.75);
    \node[model point] at (4.25,5.75) {};
    \node[above right] at (4.25,5.75) {$(17/4,23/4)$};
    \draw[model flow] (1.2,8)--(3.4,5.8);
    \node[below,align=center] at (2.5,-1.4)
    {$x_1=x_3=0$ slice; $x_4=10-x_2$\\feasible: $0\leq x_2\leq17/4$};
    \node[above] at (2.3,11.5) {objective $10+2x_2$ increases to the endpoint};
  \end{tikzpicture}
  \caption{The two zero bound coordinates reveal a one-dimensional slice.
    The multiplier certificate proves optimality in all four dimensions.}
\end{figure}

\section{Connecting multipliers, prices, and reduced costs}\label{sec:connections}
For $\max c^Tx$ with $Gx\leq b$, $x\geq0$, include bounds $-x\leq0$ in
the Lagrangian. Stationarity reads
\[
  -c+G^Ty-\mu=0,\qquad y\geq0,\quad\mu\geq0.
\]
Thus $\mu=G^Ty-c$ is the vector of decision-variable indicators.
Complementarity is simultaneously $y_i(b_i-G_{i:}x)=0$ and $\mu_jx_j=0$.
For standard equalities, use free $y$ and the same
$\mu=A^Ty-c=q$. These are the same optimality conditions expressed through
a tableau, a dual problem, and a Lagrangian.

\begin{center}
  \begin{tabular}{>{\raggedright\arraybackslash}p{3.4cm}>{\raggedright\arraybackslash}p{4.6cm}>{\raggedright\arraybackslash}p{5.1cm}}
    \toprule
    Quantity & Algebra & Interpretation\\\midrule
    RHS multiplier & $y=A_B^{-T}c_B$ & Local value of relaxing an RHS\\
    Decision indicator & $q_j=a_j^Ty-c_j$ & Marginal loss from a nonbasic activity\\
    Bound multiplier & $\mu_j=q_j$ & Nonnegative weight on $-x_j\leq0$\\
    Active-row multiplier & $\lambda_i\geq0$ & Weight of its outward normal\\\bottomrule
  \end{tabular}
\end{center}
A tight row need not have a positive price. Prices can be nonunique,
especially at degeneracy. For fixed $A,c$, the maximum value
$V(b)=\min_{A^Ty\geq c}b^Ty$ is concave in $b$ on the region of finite
feasible optima: it is a pointwise infimum of affine functions.
An old optimal $y$ gives the global upper bound
\[
  V(b+\delta b)\leq V(b)+y^T\delta b
\]
whenever the changed primal is feasible with a finite optimum.
Equality holds when the same certifying basis remains primal feasible.
At a kink, the bound remains valid even though one global derivative need
not exist. For fixed feasible set, the maximum value as a function of $c$ is
convex, because it is a supremum of linear functions $c^Tx$.
These curvature statements use \emph{maximization}; minimization reverses them.

\section{Beyond LP: what nonlinear KKT conditions establish}\label{sec:nonlinear-kkt}
For differentiable minimization with $x\in\R^n$,
\[
  \min f(x),\qquad g_i(x)\leq0,\qquad h_j(x)=0,
  \qquad L=f+\sum_i\lambda_i g_i+\sum_j\nu_jh_j,
\]
stationarity becomes
$\nabla f+\sum_i\lambda_i\nabla g_i+\sum_j\nu_j\nabla h_j=0$.
Feasibility, $\lambda\geq0$ and $\lambda_i g_i=0$ are still required.
They must be interpreted with the following distinctions.
\begin{itemize}[leftmargin=*]
  \item If $f,g_i$ are convex and $h_j$ are affine, every KKT point is a
        global minimizer. For any feasible $x$, complementarity gives
        $L(\hat x,\lambda,\nu)=f(\hat x)$ and the signs give
        $L(x,\lambda,\nu)\leq f(x)$. Convexity and stationarity make
        $\hat x$ minimize $L$, hence $f(\hat x)\leq f(x)$.
  \item In this differentiable convex setting with functions defined on
        $\R^n$, a sufficient Slater condition is a point satisfying all
        equalities and $g_i(x)<0$ for every inequality. If a minimizer exists,
        this condition guarantees KKT multipliers. More general domains use
        relative-interior qualifications; strict feasibility is sufficient,
        not necessary for every model, particularly LPs.
  \item In a general smooth problem, a local minimizer admits KKT multipliers
        under suitable regularity. For example, linear independence of the
        active inequality gradients together with the equality gradients is
        the \emph{linear independence constraint qualification} (LICQ).
        Even with this qualification, a KKT point need not be a minimum.
\end{itemize}

For a direct counterexample, minimize $f(t)=t^3$ on $[-1,1]$.
At $t=0$ both bound inequalities are inactive, their multipliers are zero,
and $f'(0)=0$: all KKT conditions hold. Yet negative nearby $t$ have smaller
value, so this point is not even a local minimum. The failure is nonconvexity,
not a missing multiplier. Likewise, $f(t)=-t^2$ has a stationary local maximum
at zero. First-order stationarity alone does not classify a point.

If functions are twice differentiable, second-order analysis examines
$\nabla^2_{xx}L$ on appropriate feasible or critical directions, not necessarily
on all of $\R^n$. For instance, minimize $u^2-v^2$ subject to $v=0$.
At $(0,0)$ the Hessian is $\operatorname{diag}(2,-2)$, but every nonzero
tangent direction is $(d,0)$ and has curvature $2d^2>0$; the constrained
point is the unique global minimum. For unconstrained problems, a
positive-definite Hessian at a stationary point suffices for a strict local
minimum; a positive-semidefinite Hessian alone does not, as $t^3$ shows.

\section{A first look at barriers and the central path}\label{sec:barrier}
Complementarity also explains the link to interior-point methods.
For the standard maximization LP, minimize $-c^Tx$ and introduce a logarithmic
barrier with parameter $\tau>0$:
\[
  \min_{Ax=b,\ x>0}\left(-c^Tx-\tau\sum_{j=1}^n\log x_j\right).
\]
Here $x>0$ means every coordinate is strictly positive; assume such a feasible
point exists and that the barrier subproblem has a minimizer.
The logarithm acts on nondimensionalized coordinates. With positive reference
scales, $\log(x_j/s_j)$ gives the same minimizers by adding a constant.
The barrier tends to $+\infty$ as a coordinate approaches zero, keeping the
subproblem's solutions in the interior relative to $Ax=b$.

Using a free equality multiplier $y$, stationarity gives
\[
  -c+A^Ty-\tau\begin{pmatrix}1/x_1\\\vdots\\1/x_n\end{pmatrix}=0.
\]
Set $q=A^Ty-c$. Then a barrier optimum satisfies
\begin{equation}\label{eq:central-path}
  Ax=b,\quad x>0,\qquad q>0,\qquad x_jq_j=\tau\quad(j=1,\ldots,n).
\end{equation}
Both primal and dual are feasible and their gap is
$b^Ty-c^Tx=x^Tq=n\tau$.
Decreasing $\tau$ therefore approaches zero complementarity gap. The family
of exact subproblem solutions is the \emph{central path}; actual algorithmic
iterates need not lie exactly on it.

Consider $\max x_1$ subject to $x_1+x_2=1$, $x\geq0$.
Here $y=\tau/x_2$ and $y-1=\tau/x_1$; eliminating $y$ gives
\[
  x_1(\tau)=\frac{1-2\tau+\sqrt{1+4\tau^2}}2,\qquad x_2=1-x_1.
\]
For every $\tau>0$, $1/2<x_1<1$; as $\tau\downarrow0$, $x$ approaches
the boundary optimum $(1,0)$. The dual slack pair stays positive until the limit.

\begin{figure}[H]
  \centering
  \begin{tikzpicture}[x=6cm,y=4cm,font=\small]
    \draw[model axis] (0,0)--(1.18,0) node[right] {$x_1$};
    \draw[model axis] (0,0)--(0,1.2) node[above] {$x_2$};
    \draw[modelblue,thick] (0,1)--(1,0);
    \node[model point,label=above left:{$\tau=1$}] at (.618034,.381966) {};
    \node[model point,label=above right:{$\tau=1/4$}] at (.809017,.190983) {};
    \node[model point,label=above right:{$\tau=1/10$}] at (.909902,.090098) {};
    \node[model point,fill=modelorange,label=below:{$x^*=(1,0)$}] at (1,0) {};
    \draw[model flow,draw=modelorange] (.67,.33)--(.76,.24);
    \node[align=left,anchor=west] at (1.35,.4)
    {$x_1+x_2=1$\\$x_jq_j=\tau$\\gap $=2\tau$};
  \end{tikzpicture}
  \caption{Computed central-path locations from the explicit formula, rather
    than solver iterates. The relative interior here is a line segment;
    decreasing the barrier parameter approaches its optimal endpoint.}
  \label{fig:central-path}
\end{figure}
Strict primal feasibility alone does not guarantee a barrier minimizer.
For example, adding a variable fixed only by $x_3\geq0$ with zero cost makes
$-\tau\log x_3$ unbounded below as $x_3\to\infty$.
Rows forcing a variable to zero also leave this simple barrier with no domain.
Such cases require reformulation or a more general initialization strategy.
Primal--dual methods solve perturbed KKT equations using Newton steps, with
step lengths controlling positivity and residuals. Their convergence and
complexity require additional analysis; the one-dimensional formula is an
illustration, not a general implementation.

\section{MATLAB checks and a study checklist}\label{sec:matlab}
The matching base MATLAB script is
\href{../../matlab/exercises/lecture_05_checks.m}{\texttt{matlab/exercises/lecture\_05\_checks.m}}.
From the repository root:
\begin{lstlisting}
run('matlab/exercises/lecture_05_checks.m')
\end{lstlisting}
The basic calculations use linear solves:
\begin{lstlisting}
AB = A(:, B);
T = AB \ A;
basic = AB \ b;
y = AB' \ c(B);
q = A' * y - c;
value = c(B)' * basic;
\end{lstlisting}
The script checks the same exact inputs as the hand solutions with absolute
tolerance $10^{-10}$. It verifies selected pivots and optimality certificates,
and enumerates all nonsingular bases in these small examples to compare
finite optima. Enumeration is a cross-check for the supplied models, not a
general solver or a test for every kind of unboundedness.
No Optimization Toolbox is required.

\begin{center}
  \begin{tabular}{>{\raggedright\arraybackslash}p{7.2cm}l}\toprule
    Hand result & Expected value\\\midrule
    Main example, second RHS change & $-50/3\leq t\leq22$\\
    Main example, decision-cost ranges & $c_1'\geq-3/2,\ c_2'\geq2,\ c_3'\leq2/5$\\
    Small model, new objective $(3,4,2)$ & $x=(8/3,10/3,0),\ z=64/3$\\
    Small model, RHS $(3,4)$ & $x=(3,0,0),\ z=6$\\
    Small model, added column & $x_1=16,\ x_6=10,\ z=42$\\
    Small model, added inequality & $x=(10/3,0,8/3),\ z=28/3$\\
    Small model, delete $x_1$ & $(x_2,x_3)=(0,6),\ z=6$\\
    Production, remove first row & $x=(4,4),\ z=20$\\
    Mixed-sign KKT certificate & $\hat x=(0,17/4,0,23/4),\ z=37/2$\\\bottomrule
  \end{tabular}
\end{center}

\subsection{Common mistakes}
\begin{center}
  \begin{tabular}{>{\raggedright\arraybackslash}p{6.3cm}>{\raggedright\arraybackslash}p{6.8cm}}\toprule
    Mistake & Correction\\\midrule
    Confusing a new coefficient with its change & Use $b_r'=b_r+t$ or $c_k'=c_k+t$\\
    Claiming unchanged basis means unchanged point & RHS changes move the basic vector\\
    Combining individual intervals independently & Test the complete vectors together\\
    Reading every slack or artificial as $+e_r$ & Check the original column and its sign\\
    Appending a raw constraint to a final tableau & Eliminate the old basic columns first\\
    Calling every basic artificial infeasible & Its optimal value must be strictly positive\\
    Treating a slack row as globally redundant & Deleting it can enlarge the feasible set\\
    Using positive multipliers on a $\geq$ normal & Flip the row to $\leq$ first\\
    Ignoring bound normals in stationarity & Include $-x_k\leq0$\\
    Claiming convexity alone guarantees KKT necessity & Check a constraint qualification\\\bottomrule
  \end{tabular}
\end{center}

\subsection{Questions to test understanding}
\begin{enumerate}[leftmargin=*]
  \item In \eqref{eq:main-lp}, what happens when $b_2'=40$?\\
        \emph{Answer:} the old basic slack is $-3$, so that basis fails. The new
        optimum is $(31/5,58/5,0)$, value $143/5$, second slack 6.
  \item What changes when $c_1'=-2$?\\
        \emph{Answer:} $t=-1<-1/2$. The old point is feasible, but its $q_4=-1/5$.
        Entering $x_4$ has a positive step and improves the new objective.
  \item For \eqref{eq:small-lp}, change both $b$ to $(3,4)$ and costs to
        $(4,-1,1)$. Why is the value not $12+(-6)+12$?\\
        \emph{Answer:} the basis passes both tests; its new point is $(3,0,0)$ and
        value 12. The separate first-order changes miss the cross term $-6$.
  \item What if a new variable has $a=0$, $\gamma=1$?\\
        \emph{Answer:} $q_w=-1$, $u=0$. The original point plus $w=\theta$ is a
        feasible improving ray, so the extended maximum is unbounded.
  \item Can a tight constraint be deleted without changing the optimum?\\
        \emph{Answer:} yes. For example, duplicate $x_1\leq1$ in $\max x_1$,
        $x_1\geq0$; deleting either copy leaves the other bound.
  \item Why do positive $\mu_1,\mu_3$ prove uniqueness in the KKT example?\\
        \emph{Answer:} any optimizer must satisfy $x_1=x_3=0$ by zero dual gap.
        The positive third-row multiplier forces that row tight; with the equality,
        these determine $x_2,x_4$ uniquely.
\end{enumerate}

\clearpage
\section{Sources and editorial remarks}\label{sec:sources}
{\small
  The primary source is E. Pawelec, Wydział MiNI PW,
  \texttt{s4\_w5.pdf}, in the external historical lecture folder.
  It has 10 PDF pages containing printed slides 1--38, with the last slide a
  closing graphic. The metadata creation date is 2019-11-03; this does not
  establish the delivery date. The folder is labelled 2020.
  Alignment with the current lecturer, syllabus, exercises, and labs is unconfirmed.

  \begin{center}
    \begin{tabular}{>{\raggedright\arraybackslash}p{4.4cm}>{\raggedright\arraybackslash}p{8.4cm}}\toprule
      Historical coverage & Corresponding development here\\\midrule
      Slides 1--8: standard form, tools, analysis types & Sections \ref{sec:basis} and \ref{sec:methods}\\
      Slides 9--14: one-parameter ranges & Sections \ref{sec:rhs}--\ref{sec:cost}; Polish exercise prompts retained\\
      Slides 15--19: vector changes & Section \ref{sec:vector}; unfinished objective calculation completed\\
      Slides 20--22: new column & Section \ref{sec:new-variable}\\
      Slides 23--27: new row & Section \ref{sec:new-constraint}; unfinished dual pivot completed\\
      Slide 28: remove a variable & Section \ref{sec:delete-variable}; feasibility phase made explicit\\
      Slides 29--33: remove a row & Section \ref{sec:delete-constraint}; useful production figures redrawn in TikZ\\
      Slides 34--37: Lagrange multipliers & Section \ref{sec:lagrange}; candidate question retained and fully certified\\\bottomrule
    \end{tabular}
  \end{center}
  Added material includes the basic-value and indicator plots, complete piecewise
  value function with proofs, data-space regions, simultaneous-change
  counterexample and cross term, structural certificates, feasible-direction
  and normal-cone diagrams, uniqueness proof, study questions, and MATLAB checks.

  Qualifications correct overly broad shortcuts: simultaneous changes are not
  inherently unpredictable; a failed basis indicator test is not always a failed
  point-optimality test; positive artificial value, not basis membership, is an
  infeasibility certificate; general nonlinear convex KKT necessity needs regularity.
  The structural example with a negative new-column component is treated as an
  algebraic LP update, without inventing a physical production interpretation.

  Two official supplementary references were consulted to cross-check terminology
  and hypotheses, with maximization signs adapted consistently:
  \begin{itemize}[leftmargin=*]
    \item \href{https://ocw.mit.edu/courses/6-251j-introduction-to-mathematical-programming-fall-2009/370982a2d3b22e090aa7a8d573bb5604_MIT6_251JF09_lec12.pdf}
          {MIT OCW 6.251J/15.081J, Fall 2009, Lecture 12: Sensitivity analysis},
          PDF pp.\ 2--6: basis tests, one-parameter updates, vector interpretation,
          new columns and rows. See Sections \ref{sec:rhs}--\ref{sec:new-constraint}.
    \item \href{https://www.seas.ucla.edu/~vandenbe/ee236b/lectures/duality.pdf}
          {L. Vandenberghe, UCLA ECE236B, Winter 2026, Chapter 5: Duality},
          sections on Slater's qualification, optimality conditions, and sensitivity.
          See Sections \ref{sec:lagrange}--\ref{sec:connections}.
  \end{itemize}
  These are supplements, not evidence of current course coverage. External
  download locations and status are recorded in the paired source inventory.

  Sections~\ref{sec:nonlinear-kkt}--\ref{sec:barrier} are added bridges to
  nonlinear and interior-point optimization, with unknown current-course
  alignment. The KKT qualifications and barrier terminology were cross-checked
  against \href{https://web.stanford.edu/~boyd/cvxbook/}{Boyd and Vandenberghe,
    \emph{Convex Optimization}}, Sections~5.2.3, 5.5.3 and~11.2.
  The cubic, equality-constrained curvature and two-variable central-path
  examples are original derivations; Figure~\ref{fig:central-path} uses their
  explicit formula. They are not covered by the existing MATLAB companion.

  Verification: the hand tableaus, endpoints, structural updates, and KKT identity
  were checked independently with exact rational arithmetic. All Lecture 5 MATLAB assertions passed in the user's normal terminal on
  2026-10-04, with reported version 26.2.0.3386108 (R2026b).
  The agent's launcher still exited silently; this is user-supplied execution
  evidence. The paired verification log records the exact command and output.
}
\end{document}
